\newpage
\section{Petri nets basics}

\subsection{Basic definitions}

A floor plan with evacuation arrows is a static spatial map: it tells \emph{where} things are. The BPMN diagram of the twelve-task insurance-claim process from the previous chapter is procedural: it tells \emph{what} the process does. What BPMN cannot answer is a third kind of question: at any given moment, \emph{where} are the running cases? No graphical artefact points at the activity each instance is sitting on. Petri nets provide exactly that, through the explicit notion of \emph{marking}.

Petri nets are named after Carl Adam Petri, who introduced them in his 1962 PhD thesis \emph{Kommunikation mit Automaten}.\footnote{Petri C.A., \emph{Kommunikation mit Automaten}, PhD thesis, 1962.} For this course they play the role of a \textbf{formal} and \textbf{abstract} business-process specification: formal because a marking and a firing rule fix the behaviour of the net, with no room for interpretation; abstract because the execution environment is disregarded, so the net carries the control logic and leaves resources, data, and platform to other concerns layered on top.

A net is built from two kinds of node. The \textbf{transition}, drawn as a square
\tikz[baseline=-0.5ex]{\node[bpmn task, rounded corners=0pt, minimum width=1.2em, minimum height=1.2em] {};},
is the \emph{active} component: an event the system can perform. Depending on the domain it may stand for an operation, a task, an evaluation, a decision, and so on: the formal model is silent on the interpretation, all transitions are graphically identical. The \textbf{place}, drawn as a circle
\tikz[baseline=-0.5ex]{\node[bpmn event, minimum size=1.2em] {};},
is the \emph{passive} one: a condition under which the system may sit, be that a state, a buffer, a repository of resources, or a memory location.

Places are passive but not empty: each may contain a number of \textbf{tokens}
\tikz[baseline=-0.5ex]{\node[bpmn event, minimum size=1.8em] (p) {}; \fill[accent] ($(p.center)+(-2pt,-1pt)$) circle (1.2pt); \fill[accent] ($(p.center)+(2pt,-1pt)$) circle (1.2pt); \fill[accent] ($(p.center)+(0,2.5pt)$) circle (1.2pt);}
that record the current state of the net. A token represents an instance of whatever the place stands for: a physical object, a piece of data, a resource, a message, a process case. It is the ``you are here'' marker of the map analogy, pinpointing one instance inside the structure the places and transitions give.

The two kinds of nodes are connected by \textbf{arcs}, which carry directed dependencies. Only two arc patterns are admitted, $p \to t$, where firing $t$ \emph{consumes} a token from $p$, and $t \to p$, where firing $t$ \emph{produces} a token in $p$. Place-to-place and transition-to-transition arcs are ruled out, which makes a Petri net a directed \textbf{bipartite} graph.

\begin{examplebox}{Pattern: sequence}
The linear chain $p_1 \to t_1 \to p_2 \to t_2 \to p_3$ with one token in $p_1$ models a strictly sequential two-step procedure: firing $t_1$ moves the token from $p_1$ to $p_2$, then $t_2$ can fire and move it to $p_3$.

\medskip
\begin{center}
\begin{tikzpicture}[node distance=0.55cm and 0.9cm]
  \node[bpmn event] (p1) {}; \fill[accent] (p1.center) circle (2pt);
  \node[bpmn task, rounded corners=0pt, right=of p1] (t1) {};
  \node[bpmn event, right=of t1] (p2) {};
  \node[bpmn task, rounded corners=0pt, right=of p2] (t2) {};
  \node[bpmn event, right=of t2] (p3) {};
  \draw[bpmn flow] (p1) -- (t1);
  \draw[bpmn flow] (t1) -- (p2);
  \draw[bpmn flow] (p2) -- (t2);
  \draw[bpmn flow] (t2) -- (p3);
  \node[below=1pt of p1, font=\scriptsize] {$p_1$};
  \node[below=1pt of t1, font=\scriptsize] {$t_1$};
  \node[below=1pt of p2, font=\scriptsize] {$p_2$};
  \node[below=1pt of t2, font=\scriptsize] {$t_2$};
  \node[below=1pt of p3, font=\scriptsize] {$p_3$};
\end{tikzpicture}
\end{center}
\end{examplebox}

\begin{examplebox}{Pattern: XOR-split}
From $p_1$ (with one token) two transitions $t_1$ and $t_2$ are both enabled: whichever fires first consumes the single token, so the other becomes disabled. The result is an exclusive choice between two downstream branches. At the Petri-net level the choice is purely non-deterministic.

\medskip
\begin{center}
\begin{tikzpicture}[node distance=0.4cm and 0.9cm]
  \node[bpmn event] (p1) {}; \fill[accent] (p1.center) circle (2pt);
  \node[bpmn task, rounded corners=0pt, above right=of p1] (t1) {};
  \node[bpmn task, rounded corners=0pt, below right=of p1] (t2) {};
  \node[bpmn event, right=of t1] (p2) {};
  \node[bpmn event, right=of t2] (p3) {};
  \draw[bpmn flow] (p1) -- (t1);
  \draw[bpmn flow] (p1) -- (t2);
  \draw[bpmn flow] (t1) -- (p2);
  \draw[bpmn flow] (t2) -- (p3);
\end{tikzpicture}
\end{center}
\end{examplebox}

\begin{examplebox}{Pattern: XOR-join}
Two input places $p_1$ and $p_2$ each feed a distinct transition ($t_1$ and $t_2$ respectively), and both transitions produce into the same place $p_3$. Whichever input place gets a token first lets the corresponding transition fire, merging the alternative paths into a common continuation.

\medskip
\begin{center}
\begin{tikzpicture}[node distance=0.4cm and 0.9cm]
  \node[bpmn event] (p1) {}; \fill[accent] (p1.center) circle (2pt);
  \node[bpmn event, below=of p1] (p2) {};
  \node[bpmn task, rounded corners=0pt, right=of p1] (t1) {};
  \node[bpmn task, rounded corners=0pt, right=of p2] (t2) {};
  \node[bpmn event, right=of t1, yshift=-0.5cm] (p3) {};
  \draw[bpmn flow] (p1) -- (t1);
  \draw[bpmn flow] (p2) -- (t2);
  \draw[bpmn flow] (t1) -- (p3);
  \draw[bpmn flow] (t2) -- (p3);
\end{tikzpicture}
\end{center}
\end{examplebox}

\begin{examplebox}{Pattern: AND-split}
From one input place $p_1$ a single transition $t_1$ produces into two output places $p_2$ and $p_3$ simultaneously. Firing $t_1$ consumes one token from $p_1$ and deposits one token in each of $p_2$ and $p_3$, so the two downstream branches become concurrently active.

\medskip
\begin{center}
\begin{tikzpicture}[node distance=0.4cm and 0.9cm]
  \node[bpmn event] (p1) {}; \fill[accent] (p1.center) circle (2pt);
  \node[bpmn task, rounded corners=0pt, right=of p1] (t1) {};
  \node[bpmn event, above right=of t1] (p2) {};
  \node[bpmn event, below right=of t1] (p3) {};
  \draw[bpmn flow] (p1) -- (t1);
  \draw[bpmn flow] (t1) -- (p2);
  \draw[bpmn flow] (t1) -- (p3);
\end{tikzpicture}
\end{center}
\end{examplebox}

\begin{examplebox}{Pattern: AND-join}
Two input places $p_1$ and $p_2$ feed a single transition $t_1$ that produces into $p_3$. The transition is enabled only when \emph{both} input places hold a token; firing consumes one token from each and produces one in $p_3$. This is the synchronisation construct: it waits for both branches before continuing.

\medskip
\begin{center}
\begin{tikzpicture}[node distance=0.4cm and 0.9cm]
  \node[bpmn event] (p1) {}; \fill[accent] (p1.center) circle (2pt);
  \node[bpmn event, below=of p1] (p2) {}; \fill[accent] (p2.center) circle (2pt);
  \node[bpmn task, rounded corners=0pt, right=of p1, yshift=-0.5cm] (t1) {};
  \node[bpmn event, right=of t1] (p3) {};
  \draw[bpmn flow] (p1) -- (t1);
  \draw[bpmn flow] (p2) -- (t1);
  \draw[bpmn flow] (t1) -- (p3);
\end{tikzpicture}
\end{center}
\end{examplebox}

\begin{examplebox}{Pattern: OR-split (inclusive choice)}
From $p_1$ (with one token) three transitions are simultaneously available: $t_2$ produces into $p_2$ only, $t_3$ produces into $p_3$ only, and $t_1$, an AND-split in disguise, produces into both. The net thus expresses the inclusive disjunction ``$p_2$ or $p_3$ or both''.

\medskip
\begin{center}
\begin{tikzpicture}[node distance=0.35cm and 0.9cm]
  \node[bpmn event] (p1) {}; \fill[accent] (p1.center) circle (2pt);
  \node[bpmn task, rounded corners=0pt, above right=of p1] (t2) {};
  \node[bpmn task, rounded corners=0pt, right=of p1] (t1) {};
  \node[bpmn task, rounded corners=0pt, below right=of p1] (t3) {};
  \node[bpmn event, right=of t2] (p2) {};
  \node[bpmn event, right=of t3] (p3) {};
  \draw[bpmn flow] (p1) -- (t2);
  \draw[bpmn flow] (p1) -- (t1);
  \draw[bpmn flow] (p1) -- (t3);
  \draw[bpmn flow] (t2) -- (p2);
  \draw[bpmn flow] (t3) -- (p3);
  \draw[bpmn flow] (t1) -- (p2);
  \draw[bpmn flow] (t1) -- (p3);
\end{tikzpicture}
\end{center}
\end{examplebox}

\begin{examplebox}{Cappuccino}
Tokens also carry the intuition of \emph{consumable resources}: the inputs of a transition are what it needs, the outputs what it yields. Two input places, \textsc{milk} (holding two tokens) and \textsc{coffee} (holding three tokens), both feed a single transition; firing it consumes \emph{one} milk token and \emph{one} coffee token and produces one \textsc{cappuccino} token.

\medskip
\begin{center}
\begin{tikzpicture}[node distance=0.4cm and 0.9cm]
  \node[bpmn event, minimum size=2.6em] (milk) {};
  \fill[accent] ($(milk.center)+(-3pt,0)$) circle (1.5pt);
  \fill[accent] ($(milk.center)+(3pt,0)$) circle (1.5pt);
  \node[bpmn event, minimum size=2.6em, below=of milk] (coffee) {};
  \fill[accent] ($(coffee.center)+(-4pt,0)$) circle (1.5pt);
  \fill[accent] ($(coffee.center)+(0,0)$) circle (1.5pt);
  \fill[accent] ($(coffee.center)+(4pt,0)$) circle (1.5pt);
  \node[bpmn task, rounded corners=0pt, right=of milk, yshift=-0.5cm] (t) {};
  \node[bpmn event, minimum size=2.4em, right=of t] (capp) {};
  \fill[accent] (capp.center) circle (1.5pt);
  \draw[bpmn flow] (milk) -- (t);
  \draw[bpmn flow] (coffee) -- (t);
  \draw[bpmn flow] (t) -- (capp);
  \node[left=1pt of milk, font=\scriptsize] {milk};
  \node[left=1pt of coffee, font=\scriptsize] {coffee};
  \node[right=1pt of capp, font=\scriptsize] {cappuccino};
\end{tikzpicture}
\end{center}
\end{examplebox}

\begin{definitionbox}{Petri net}
A \textbf{Petri net} is a tuple $(P, T, F, M_0)$ where:
\begin{itemize}
  \item $P$ is a finite set of \emph{places};
  \item $T$ is a finite set of \emph{transitions}, with $P \cap T = \emptyset$;
  \item $F \subseteq (P \times T) \cup (T \times P)$ is the \emph{flow relation} (the arcs);
  \item $M_0: P \to \mathbb{N}$ is the \emph{initial marking}.
\end{itemize}
\end{definitionbox}

\begin{definitionbox}{Pre-set and post-set}
The \textbf{pre-set} $\bullet t$ of a transition is the set of its \emph{input places}, those the tokens needed to fire are drawn from; the \textbf{post-set} $t \bullet$ is the set of its \emph{output places}, those the produced tokens land in. The dot sits on the side the tokens come from. Both read off the flow relation, and the same reading works for any node $x \in P \cup T$:
\[
  \bullet x \;=\; \{ y \mid (y,x) \in F \}, \qquad x \bullet \;=\; \{ y \mid (x,y) \in F \}.
\]
For a place $p$ the roles swap: $\bullet p$ holds the transitions tokens in $p$ \emph{can come from}, $p \bullet$ those they \emph{can go to}. The bipartite constraint keeps the two apart, forcing $\bullet t, t \bullet \subseteq P$ and $\bullet p, p \bullet \subseteq T$.
\end{definitionbox}

\subsection{Token game}

\begin{definitionbox}{Enabling and firing}
A transition $t$ is \textbf{enabled} if each of its input places contains at least one token. When an enabled transition \textbf{fires}, it \emph{consumes} one token from each input place and \emph{produces} one token into each output place.
\end{definitionbox}

\begin{notebox}{Remarks on firing}
\begin{itemize}
  \item \textbf{Atomicity}: a firing is an atomic action: the net never lingers in a half-fired state.
  \item \textbf{Interleaving semantics}: multiple transitions may be enabled at the same time, but only one fires at a time. Concurrency is rendered as the set of all possible interleavings of independent firings.
  \item \textbf{Static structure, variable population}: the net itself ($P$, $T$, $F$) is fixed, but the total number of tokens may grow or shrink over time. Firing a transition $t$ adds $|t\bullet| - |\bullet t|$ tokens to the system, which can be positive, zero, or negative.
\end{itemize}
\end{notebox}

\begin{examplebox}{Token game (start)}
The net has four places $p_1, p_2, p_3, p_4$ and four transitions $t_1, t_2, t_3, t_4$, arranged in a grid with $p_3$ at the centre: $t_1$ consumes from $p_1$ and produces into \emph{both} $p_2$ and $p_3$; $t_3$ consumes from \emph{both} $p_2$ and $p_3$ and produces into $p_4$; $t_2$ moves a token from $p_3$ back to $p_1$; $t_4$ moves a token from $p_4$ back to $p_3$. The initial marking puts a single token in $p_1$, written $M_0 = p_1$, and only $t_1$ is enabled. The table tracks one run as a sequence alternating markings and fired transitions.

\medskip
\begin{center}
\begin{minipage}{0.46\linewidth}
\centering
\begin{tikzpicture}[node distance=0.6cm and 1.0cm]
  \node[bpmn event, minimum size=1.6em] (p1) {}; \fill[accent] (p1.center) circle (2pt);
  \node[bpmn task, rounded corners=0pt, minimum size=1.6em, right=of p1] (t1) {};
  \node[bpmn event, minimum size=1.6em, right=of t1] (p2) {};
  \node[bpmn task, rounded corners=0pt, minimum size=1.6em, below=of p1] (t2) {};
  \node[bpmn event, minimum size=1.6em, below=of t1] (p3) {};
  \node[bpmn task, rounded corners=0pt, minimum size=1.6em, below=of p2] (t3) {};
  \node[bpmn task, rounded corners=0pt, minimum size=1.6em, below=of p3] (t4) {};
  \node[bpmn event, minimum size=1.6em, below=of t3] (p4) {};
  \draw[bpmn flow] (p1) -- (t1);
  \draw[bpmn flow] (t1) -- (p2);
  \draw[bpmn flow] (t1) -- (p3);
  \draw[bpmn flow] (p2) -- (t3);
  \draw[bpmn flow] (p3) -- (t3);
  \draw[bpmn flow] (p3) -- (t2);
  \draw[bpmn flow] (t2) -- (p1);
  \draw[bpmn flow] (t3) -- (p4);
  \draw[bpmn flow] (p4) -- (t4);
  \draw[bpmn flow] (t4) -- (p3);
  \node[below left=0pt of p1, font=\scriptsize] {$p_1$};
  \node[below left=0pt of t1, font=\scriptsize] {$t_1$};
  \node[below left=0pt of p2, font=\scriptsize] {$p_2$};
  \node[below left=0pt of t2, font=\scriptsize] {$t_2$};
  \node[below left=0pt of p3, font=\scriptsize] {$p_3$};
  \node[below left=0pt of t3, font=\scriptsize] {$t_3$};
  \node[below left=0pt of t4, font=\scriptsize] {$t_4$};
  \node[below left=0pt of p4, font=\scriptsize] {$p_4$};
\end{tikzpicture}
\end{minipage}%
\hfill
\begin{minipage}{0.5\linewidth}
\centering\small
\begin{tabular}{lll}
\toprule
\textbf{Fired} & \textbf{New marking} & \textbf{Enabled} \\
\midrule
--- & $p_1$ & $t_1$ \\
$t_1$ & $p_2+p_3$ & $t_2$, $t_3$ \\
$t_2$ & $p_1+p_2$ & $t_1$ \\
$t_1$ & $2\cdot p_2+p_3$ & $t_2$, $t_3$ \\
$t_3$ & $p_2+p_4$ & $t_4$ \\
$t_4$ & $p_2+p_3$ & $t_2$, $t_3$ \\
\bottomrule
\end{tabular}
\end{minipage}
\end{center}

\medskip
The choice between $t_2$ and $t_3$ after $t_1$ is non-deterministic. The marking $2\cdot p_2+p_3$ puts two tokens in $p_2$: a marking is in general a \emph{multiset} over $P$, not a subset; and firing $t_3$ there consumes one token from $p_2$ \emph{and} one from $p_3$, leaving $p_2+p_4$. The run is cyclic: $p_2+p_3$ reappears and the same choice opens again. Taking $t_3$ at the first opportunity gives the second run, $p_2+p_3 \;/\; t_3 \;/\; p_4 \;/\; t_4 \;/\; p_3 \;/\; t_2 \;/\; p_1$: back to the initial marking.
\end{examplebox}

\subsection{Workflow nets}

Workflow nets are a restricted subclass of Petri nets, tailored to business-process modelling: the structural restrictions match the lifecycle of a process instance, a single start, a single end, and a path connecting them. The semantics stays the one of the token game, so every modelled process keeps a precise meaning, and transitions and arcs may carry decorations (split/join roles, triggers) that BPMN-style readers interpret directly.

\medskip
\begin{center}
\begin{tikzpicture}[node distance=0.4cm and 1.1cm]
  \node[bpmn event, minimum size=1.6em] (s) {};
  \node[bpmn task, fill=rowgray, rounded corners=6pt, minimum width=2.6em, minimum height=2em, right=of s] (b) {WFN};
  \node[bpmn event, minimum size=1.6em, right=of b] (e) {};
  \draw[bpmn flow] (s) -- (b);
  \draw[bpmn flow] (b) -- (e);
  \node[below=2pt of s, font=\scriptsize] {start};
  \node[below=2pt of e, font=\scriptsize] {end};
\end{tikzpicture}
\end{center}

A token entering at \emph{start} represents one new process instance; the token reaching \emph{end} represents that instance being completed.

\begin{definitionbox}{Workflow net}
Being a workflow net is a property of the \emph{structure} alone, so the initial marking plays no part and the net is read as the triple $(P, T, F)$. It is a \textbf{workflow net} (WFN) iff:
\begin{enumerate}
  \item there is a distinguished \emph{initial place} $i \in P$ with $\bullet i = \emptyset$ (no incoming arcs);
  \item there is a distinguished \emph{final place} $o \in P$ with $o \bullet = \emptyset$ (no outgoing arcs);
  \item every other place and every transition belongs to at least one path from $i$ to $o$.
\end{enumerate}
\end{definitionbox}

\begin{notebox}{Basic property: uniqueness of the initial node}
\textbf{Lemma.} In a workflow net there is a \emph{unique} node with no incoming arc.

\medskip

\textit{Proof.} Let $i$ be the initial place and $o$ the final one. Suppose, for contradiction, that some other node $v$ has $\bullet v = \emptyset$. By condition~(3) of the workflow-net definition, $v$ lies on some path from $i$ to $o$; since $v$ has no incoming arc, it must be the \emph{first} node of that path. The first node of a path from $i$ is $i$ itself, hence $v = i$. \qed
\end{notebox}

\begin{notebox}{Basic property: uniqueness of the final node}
\textbf{Lemma.} In a workflow net there is a \emph{unique} node with no outgoing arc: namely the final place $o$. The proof is analogous to the previous lemma, swapping pre-set for post-set and reading the path from $o$ backwards.
\end{notebox}

Condition (3) is what rules out orphans: a node on no path from $i$ to $o$ would sit outside every case the net can run.

\begin{examplebox}{Workflow net: example}
An order-handling process modelled as a workflow net. The token enters at $i$; \textsc{receive\_order} produces into $p_2$; \textsc{process\_order} is an AND-split that produces tokens into $p_3$ and $p_4$ simultaneously; the two branches run \textsc{pack\_goods} (producing into $p_5$) and \textsc{update\_inventory} (producing into $p_6$); the AND-join \textsc{complete\_order} synchronises $p_5$ and $p_6$ and produces into $p_7$; finally \textsc{send\_goods} delivers the token to $o$. Eight places, six transitions, single source, single sink, every node on some path from $i$ to $o$.

\medskip
\begin{center}
\begin{tikzpicture}
  \node[bpmn event, minimum size=1.5em] (i) at (0,0) {};
  \node[bpmn task, rounded corners=0pt, font=\scriptsize\sffamily] (rec) at (1.7,0) {receive order};
  \node[bpmn event, minimum size=1.5em] (p2) at (3.4,0) {};
  \node[bpmn task, rounded corners=0pt, font=\scriptsize\sffamily] (proc) at (5.1,0) {process order};
  \node[bpmn event, minimum size=1.5em] (p3) at (6.8,0.9) {};
  \node[bpmn event, minimum size=1.5em] (p4) at (6.8,-0.9) {};
  \node[bpmn task, rounded corners=0pt, font=\scriptsize\sffamily] (pack) at (8.5,0.9) {pack goods};
  \node[bpmn task, rounded corners=0pt, font=\scriptsize\sffamily] (upd) at (8.5,-0.9) {update inventory};
  \node[bpmn event, minimum size=1.5em] (p5) at (10.2,0.9) {};
  \node[bpmn event, minimum size=1.5em] (p6) at (10.2,-0.9) {};
  \node[bpmn task, rounded corners=0pt, font=\scriptsize\sffamily] (comp) at (11.9,0) {complete order};
  \node[bpmn event, minimum size=1.5em] (p7) at (13.6,0) {};
  \node[bpmn task, rounded corners=0pt, font=\scriptsize\sffamily] (send) at (15.1,0) {send goods};
  \node[bpmn event, minimum size=1.5em] (o) at (16.6,0) {};
  \draw[bpmn flow] (i) -- (rec);
  \draw[bpmn flow] (rec) -- (p2);
  \draw[bpmn flow] (p2) -- (proc);
  \draw[bpmn flow] (proc) -- (p3);
  \draw[bpmn flow] (proc) -- (p4);
  \draw[bpmn flow] (p3) -- (pack);
  \draw[bpmn flow] (p4) -- (upd);
  \draw[bpmn flow] (pack) -- (p5);
  \draw[bpmn flow] (upd) -- (p6);
  \draw[bpmn flow] (p5) -- (comp);
  \draw[bpmn flow] (p6) -- (comp);
  \draw[bpmn flow] (comp) -- (p7);
  \draw[bpmn flow] (p7) -- (send);
  \draw[bpmn flow] (send) -- (o);
  \node[below=1pt of i, font=\scriptsize] {$i$};
  \node[below=1pt of p2, font=\scriptsize] {$p_2$};
  \node[below=1pt of p3, font=\scriptsize] {$p_3$};
  \node[below=1pt of p4, font=\scriptsize] {$p_4$};
  \node[below=1pt of p5, font=\scriptsize] {$p_5$};
  \node[below=1pt of p6, font=\scriptsize] {$p_6$};
  \node[below=1pt of p7, font=\scriptsize] {$p_7$};
  \node[below=1pt of o, font=\scriptsize] {$o$};
\end{tikzpicture}
\end{center}
\end{examplebox}

\begin{examplebox}{Question time: four variants}
Four small nets test the definition.
\begin{enumerate}
  \item Places $i, p_2, o$; arcs $i \to t_1 \to p_2$, $p_2 \to t_3 \to o$, $p_2 \to t_2 \to o$. \textbf{WFN}: $\bullet i = \emptyset$, $o\bullet = \emptyset$, and every node lies on a path such as $i \to t_1 \to p_2 \to t_2 \to o$.
  \item Add a back-arc $t_2 \to i$. \textbf{Not a WFN}: $i$ now has an incoming arc, so there is no proper initial place, however well-connected the net is.
  \item Extend the first net instead with a side branch $p_2 \to t_2 \to p_4 \to t_4 \to p_5$ ending in the dangling place $p_5$. \textbf{Not a WFN}: $t_2$, $p_4$, $t_4$, $p_5$ lie on no path from $i$ to $o$: a dead end violates condition~(3).
  \item Close the dead end with a new transition $t_5$ from $p_5$ back to $p_2$. \textbf{WFN again}: every loop node now reaches $o$, e.g.\ via $i \to t_1 \to p_2 \to t_2 \to p_4 \to t_4 \to p_5 \to t_5 \to p_2 \to t_3 \to o$.
\end{enumerate}

\medskip
% Verified 2026-07-10: two structure nets (unmarked). Variant 1: places i,p_2,o;
% transitions t_1,t_2,t_3; arcs i->t_1, t_1->p_2, p_2->t_3, t_3->o, p_2->t_2, t_2->o.
% Unique source i (no in-arc), unique sink o (no out-arc); every node on an i->o path
% (i->t_1->p_2->t_3->o and i->t_1->p_2->t_2->o) => WFN. Variant 3: variant 1's forward
% path i->t_1->p_2->t_3->o plus side branch p_2->t_2->p_4->t_4->p_5 with p_5 dangling;
% t_2,p_4,t_4,p_5 reach no o => not a WFN (condition 3). Only arcs listed in the tex.
% Coords chosen to fit A4, no resizebox.
\begin{center}
\begin{tikzpicture}[node distance=0.7cm and 0.9cm]
  % --- Variant 1 (WFN) ---
  \node[font=\scriptsize\sffamily] at (1.6,1.3) {Variant 1: WFN};
  \node[bpmn event, minimum size=1.5em] (v1i) at (0,0) {};
  \node[bpmn task, rounded corners=0pt, minimum size=1.5em] (v1t1) at (1.1,0) {};
  \node[bpmn event, minimum size=1.5em] (v1p2) at (2.2,0) {};
  \node[bpmn task, rounded corners=0pt, minimum size=1.5em] (v1t3) at (3.3,0.6) {};
  \node[bpmn task, rounded corners=0pt, minimum size=1.5em] (v1t2) at (3.3,-0.6) {};
  \node[bpmn event, minimum size=1.5em] (v1o) at (4.4,0) {};
  \draw[bpmn flow] (v1i) -- (v1t1);
  \draw[bpmn flow] (v1t1) -- (v1p2);
  \draw[bpmn flow] (v1p2) -- (v1t3);
  \draw[bpmn flow] (v1t3) -- (v1o);
  \draw[bpmn flow] (v1p2) -- (v1t2);
  \draw[bpmn flow] (v1t2) -- (v1o);
  \node[below=1pt of v1i, font=\scriptsize] {$i$};
  \node[above=1pt of v1t1, font=\scriptsize] {$t_1$};
  \node[below=1pt of v1p2, font=\scriptsize] {$p_2$};
  \node[above=1pt of v1t3, font=\scriptsize] {$t_3$};
  \node[below=1pt of v1t2, font=\scriptsize] {$t_2$};
  \node[below=1pt of v1o, font=\scriptsize] {$o$};
  % --- Variant 3 (not a WFN: dead end) ---
  \node[font=\scriptsize\sffamily] at (8.3,1.3) {Variant 3: dead end (not a WFN)};
  \node[bpmn event, minimum size=1.5em] (v3i) at (6.4,0) {};
  \node[bpmn task, rounded corners=0pt, minimum size=1.5em] (v3t1) at (7.5,0) {};
  \node[bpmn event, minimum size=1.5em] (v3p2) at (8.6,0) {};
  \node[bpmn task, rounded corners=0pt, minimum size=1.5em] (v3t3) at (9.7,0) {};
  \node[bpmn event, minimum size=1.5em] (v3o) at (10.8,0) {};
  \node[bpmn task, rounded corners=0pt, minimum size=1.5em] (v3t2) at (8.6,-1.1) {};
  \node[bpmn event, minimum size=1.5em] (v3p4) at (9.7,-1.1) {};
  \node[bpmn task, rounded corners=0pt, minimum size=1.5em] (v3t4) at (10.8,-1.1) {};
  \node[bpmn event, minimum size=1.5em] (v3p5) at (11.9,-1.1) {};
  \draw[bpmn flow] (v3i) -- (v3t1);
  \draw[bpmn flow] (v3t1) -- (v3p2);
  \draw[bpmn flow] (v3p2) -- (v3t3);
  \draw[bpmn flow] (v3t3) -- (v3o);
  \draw[bpmn flow] (v3p2) -- (v3t2);
  \draw[bpmn flow] (v3t2) -- (v3p4);
  \draw[bpmn flow] (v3p4) -- (v3t4);
  \draw[bpmn flow] (v3t4) -- (v3p5);
  \node[below=1pt of v3i, font=\scriptsize] {$i$};
  \node[above=1pt of v3t1, font=\scriptsize] {$t_1$};
  \node[above=1pt of v3p2, font=\scriptsize] {$p_2$};
  \node[above=1pt of v3t3, font=\scriptsize] {$t_3$};
  \node[above=1pt of v3o, font=\scriptsize] {$o$};
  \node[below=1pt of v3t2, font=\scriptsize] {$t_2$};
  \node[below=1pt of v3p4, font=\scriptsize] {$p_4$};
  \node[below=1pt of v3t4, font=\scriptsize] {$t_4$};
  \node[below=1pt of v3p5, font=\scriptsize] {$p_5$};
\end{tikzpicture}
\end{center}
\end{examplebox}

\begin{examplebox}{Question time: is this a WFN?}
A larger candidate: thirteen places and twelve transitions (numbering, with gaps, follows the original drawing), wired left-to-right with several branches. Is it a workflow net?

\medskip
\begin{center}
\begin{tikzpicture}
  \node[bpmn event, minimum size=1.5em] (i)   at (0,0) {};
  \node[bpmn task, rounded corners=0pt, minimum size=1.5em] (t1)  at (2,0) {};
  \node[bpmn event, minimum size=1.5em] (p8)  at (4,0) {};
  \node[bpmn task, rounded corners=0pt, minimum size=1.5em] (t7)  at (6,0) {};
  \node[bpmn event, minimum size=1.5em] (p10) at (7.6,0) {};
  \node[bpmn task, rounded corners=0pt, minimum size=1.5em] (t9)  at (9.2,0) {};
  \node[bpmn event, minimum size=1.5em] (p13) at (11.2,0) {};
  \node[bpmn task, rounded corners=0pt, minimum size=1.5em] (t6)  at (0,-1.5) {};
  \node[bpmn event, minimum size=1.5em] (p7)  at (1.8,-1.4) {};
  \node[bpmn event, minimum size=1.5em] (p2)  at (3.9,-1.4) {};
  \node[bpmn event, minimum size=1.5em] (p9)  at (6.5,-1.4) {};
  \node[bpmn event, minimum size=1.5em] (p11) at (8.6,-1.6) {};
  \node[bpmn event, minimum size=1.5em] (p12) at (10.2,-1.4) {};
  \node[bpmn task, rounded corners=0pt, minimum size=1.5em] (t10) at (12.4,-1.5) {};
  \node[bpmn task, rounded corners=0pt, minimum size=1.5em] (t3)  at (3.2,-2.7) {};
  \node[bpmn task, rounded corners=0pt, minimum size=1.5em] (t8)  at (5.6,-2.5) {};
  \node[bpmn task, rounded corners=0pt, minimum size=1.5em] (t15) at (6.9,-2.8) {};
  \node[bpmn task, rounded corners=0pt, minimum size=1.5em] (t11) at (10.2,-2.7) {};
  \node[bpmn event, minimum size=1.5em] (p6)  at (0,-3.0) {};
  \node[bpmn task, rounded corners=0pt, minimum size=1.5em] (t12) at (7.9,-3.6) {};
  \node[bpmn event, minimum size=1.5em] (p14) at (10.7,-3.7) {};
  \node[bpmn task, rounded corners=0pt, minimum size=1.5em] (t14) at (2.4,-4.7) {};
  \node[bpmn event, minimum size=1.5em] (p15) at (5.4,-4.6) {};
  \node[bpmn task, rounded corners=0pt, minimum size=1.5em] (t16) at (8.2,-4.8) {};
  \node[bpmn event, minimum size=1.5em] (o)   at (12.6,-4.6) {};
  \draw[bpmn flow] (i) -- (t1);
  \draw[bpmn flow] (i) -- (t6);
  \draw[bpmn flow] (t1) -- (p8);
  \draw[bpmn flow] (t1) -- (p2);
  \draw[bpmn flow] (p8) -- (t7);
  \draw[bpmn flow] (t7) -- (p10);
  \draw[bpmn flow] (t7) -- (p9);
  \draw[bpmn flow] (p10) -- (t9);
  \draw[bpmn flow] (t9) -- (p13);
  \draw[bpmn flow] (t9) -- (p11);
  \draw[bpmn flow] (t9) -- (p12);
  \draw[bpmn flow] (p12) -- (t11);
  \draw[bpmn flow] (t11) -- (p11);
  \draw[bpmn flow] (p11) -- (t15);
  \draw[bpmn flow] (p11) -- (t12);
  \draw[bpmn flow] (t12) -- (p14);
  \draw[bpmn flow] (p13) -- (t10);
  \draw[bpmn flow] (p14) -- (t10);
  \draw[bpmn flow] (t10) -- (o);
  \draw[bpmn flow] (p9) -- (t8);
  \draw[bpmn flow] (p2) -- (t8);
  \draw[bpmn flow] (p2) -- (t3);
  \draw[bpmn flow] (p7) -- (t3);
  \draw[bpmn flow] (t3) -- (p15);
  \draw[bpmn flow] (t6) -- (p7);
  \draw[bpmn flow] (t6) -- (p6);
  \draw[bpmn flow] (p6) -- (t14);
  \draw[bpmn flow] (p7) -- (t14);
  \draw[bpmn flow] (t14) -- (p15);
  \draw[bpmn flow] (t8) -- (p15);
  \draw[bpmn flow] (t15) -- (p15);
  \draw[bpmn flow] (p15) -- (t16);
  \draw[bpmn flow] (t16) -- (o);
  \node[below=1pt of i, font=\scriptsize] {$i$};
  \node[below=1pt of t1, font=\scriptsize] {$t_1$};
  \node[below=1pt of p8, font=\scriptsize] {$p_8$};
  \node[below=1pt of t7, font=\scriptsize] {$t_7$};
  \node[below=1pt of p10, font=\scriptsize] {$p_{10}$};
  \node[below=1pt of t9, font=\scriptsize] {$t_9$};
  \node[below=1pt of p13, font=\scriptsize] {$p_{13}$};
  \node[below=1pt of t6, font=\scriptsize] {$t_6$};
  \node[above=1pt of p7, font=\scriptsize] {$p_7$};
  \node[above=1pt of p2, font=\scriptsize] {$p_2$};
  \node[above=1pt of p9, font=\scriptsize] {$p_9$};
  \node[left=1pt of p11, font=\scriptsize] {$p_{11}$};
  \node[above=1pt of p12, font=\scriptsize] {$p_{12}$};
  \node[right=1pt of t10, font=\scriptsize] {$t_{10}$};
  \node[below=1pt of t3, font=\scriptsize] {$t_3$};
  \node[left=1pt of t8, font=\scriptsize] {$t_8$};
  \node[below=1pt of t15, font=\scriptsize] {$t_{15}$};
  \node[below=1pt of t11, font=\scriptsize] {$t_{11}$};
  \node[below=1pt of p6, font=\scriptsize] {$p_6$};
  \node[below=1pt of t12, font=\scriptsize] {$t_{12}$};
  \node[below=1pt of p14, font=\scriptsize] {$p_{14}$};
  \node[below=1pt of t14, font=\scriptsize] {$t_{14}$};
  \node[below=1pt of p15, font=\scriptsize] {$p_{15}$};
  \node[below=1pt of t16, font=\scriptsize] {$t_{16}$};
  \node[below=1pt of o, font=\scriptsize] {$o$};
\end{tikzpicture}
\end{center}

$i$ has no incoming arc and $o$ no outgoing one, so conditions (1) and (2) hold. Condition (3) asks that every node lie on \emph{some} path from $i$ to $o$, not that one path cross them all. Most nodes are covered at a glance, e.g.\ by $i \to t_1 \to p_8 \to t_7 \to p_{10} \to t_9 \to p_{13} \to t_{10} \to o$. The delicate spot is $t_3$: without the arc $t_3 \to p_{15}$, no path through $t_3$ (or reaching it via $p_2$ or $p_7$) could continue to $o$, and the net would \emph{not} be a workflow net. With that arc in place (drawn in the figure), $i \to t_6 \to p_7 \to t_3 \to p_{15} \to t_{16} \to o$ covers the critical nodes, and the verdict is \textbf{yes}.
\end{examplebox}

\subsection{WoPeD and its notation}

The course's reference tool for designing and analysing workflow nets is \textbf{WoPeD}, the Workflow Petri Net Designer\footnote{\url{http://woped.dhbw-karlsruhe.de/woped/}}: a free desktop application offering a graphical editor, a transition-properties dialog, and automated analyses such as coverability-graph generation. Two of its features extend the plain-net notation: decorated transitions (\emph{syntax sugar}) and hierarchical sub-processes.

\begin{notebox}{Transition decorations in WoPeD}
WoPeD's transition-properties dialog lets a transition carry a \emph{branching role}: \textsc{None}, \textsc{AND-split}, \textsc{XOR-split}, \textsc{AND-join}, \textsc{XOR-join}, and the mixed forms \textsc{XOR-join-split}, \textsc{AND-join-XOR-split}, \textsc{AND-join-split}, \textsc{XOR-join-AND-split}. The decoration is purely \emph{syntax sugar}: each decorated transition stands for a sub-net of plain Petri-net transitions and places.\footnote{\url{http://woped.dhbw-karlsruhe.de/woped/}}
\end{notebox}

\textbf{Sugar for split.} A single decorated transition $t_1$ with two output places can stand for two different patterns depending on the split type:

\begin{itemize}
  \item \textbf{AND-split}. $t_1$ behaves like a plain transition with two outgoing arcs: firing produces tokens in both $p_2$ and $p_3$ simultaneously. The unfolding is that same transition.
  \item \textbf{XOR-split}. $t_1$ stands for two distinct transitions $t_{1a}$ and $t_{1b}$ that share the input place $p_1$ but produce into different outputs ($t_{1a} \to p_2$, $t_{1b} \to p_3$). Either fires, never both.
\end{itemize}

\textbf{Sugar for join.} Symmetric to the split sugar:

\begin{itemize}
  \item \textbf{AND-join}. A decorated transition $t_1$ with two input places $p_1$ and $p_2$ stands for a plain transition with two incoming arcs: firing consumes one token from each and produces into $p_3$.
  \item \textbf{XOR-join}. $t_1$ stands for two distinct transitions $t_{1a}$ and $t_{1b}$: $t_{1a}$ consumes from $p_1$ and produces into $p_3$, $t_{1b}$ consumes from $p_2$ and produces into $p_3$. Either of the two input places, when marked, lets the corresponding transition fire.
\end{itemize}

\textbf{Mixed sugar.} The two kinds of sugar combine on the same transition. A four-port decorated $t_1$ with input places $p_1, p_2$ and output places $p_3, p_4$, an \textsc{AND-join} on the input side and an \textsc{XOR-split} on the output, admits two equivalent expansions:

\begin{itemize}
  \item with an \emph{intermediate place}: $t_{1a}$ consumes $\{p_1, p_2\}$ and produces into a fresh place $pt_1$; from there $t_{1b}$ produces into $p_3$ and $t_{1c}$ into $p_4$;
  \item without: two transitions, each consuming $\{p_1, p_2\}$ and producing into one of $p_3$ or $p_4$.
\end{itemize}

\begin{notebox}{A personal remark on syntax sugar}
The decorations carry three caveats:
\begin{itemize}
  \item the chosen glyphs for AND-join and XOR-join are visually too similar (a small vertical bar versus a hatched bar) and can be confused at first reading;
  \item the same glyph placed on the input side or on the output side of a transition denotes \emph{different} patterns (join versus split), so position matters;
  \item the AND decoration is in fact \emph{redundant}: a plain transition with multiple input/output arcs already behaves as an AND-join / AND-split, so the dedicated AND-bar adds nothing semantically.
\end{itemize}
\end{notebox}

\begin{notebox}{Why the decoration sits on the transition}
The decoration is a back-port from BPMN, where the corresponding split/join roles are filled by dedicated \emph{gateway} nodes drawn as diamonds: a $+$ diamond for AND, a $\times$ diamond for XOR. In Petri nets the gateway is absorbed into the transition, but the conceptual role is the same.
\end{notebox}

\begin{notebox}{Convention adopted in the course}
To avoid any source of confusion, going forward the AND-join and AND-split decorations are \textbf{abandoned}: AND-patterns are always drawn as plain transitions with multiple input/output arcs. Only the XOR-join and XOR-split decorations are retained as legitimate syntax sugar, because their unfolding is genuinely non-trivial.
\end{notebox}

\textbf{Sub-processes.} WoPeD's second notational extension is hierarchy.

\begin{examplebox}{WoPeD: sub-process}
WoPeD supports hierarchical modelling: a transition of a parent workflow net (drawn as a \emph{double square}) can be opened as a separate workflow net of its own. A transition is therefore not necessarily atomic: it can be \emph{implemented} by another WFN, opaque from the outside.\footnote{\url{http://woped.dhbw-karlsruhe.de/woped/}}
\end{examplebox}

\begin{notebox}{Hierarchical structuring}
The uniqueness of the entry and exit points is exactly what makes hierarchical structuring possible. A transition labelled \textsc{business process A} can be \emph{realised} by an entire workflow net $A$: the transition's single input arc becomes the entry into $A$'s initial place, its single output arc the exit from $A$'s final place. From outside it looks like a single atomic step; inside it is a workflow net.
\end{notebox}

\subsection{Some patterns}

The visual-notation chapter catalogued five control-flow patterns on the insurance-claim diagram: sequence, parallel, choice, merge, iteration. As workflow nets they are the fragments already built above: sequencing is the linear chain $i \to A \to p \to B \to o$, where the middle place holds the token until $A$ has produced into it; parallelism is an AND-split feeding two branches that an AND-join synchronises; selection fuses the XOR-split and XOR-join, which BPMN keeps apart as choice and merge. Two constructs remain, and neither survives the translation to BPMN unchanged: a choice the net makes \emph{implicitly}, and the capacity constraints, which BPMN-style control flow cannot express at all.

\begin{center}
\begin{tikzpicture}[node distance=0.4cm and 0.9cm]
  \node[bpmn event] (i) {};
  \node[bpmn task, rounded corners=0pt, right=of i] (A) {$A$};
  \node[bpmn event, right=of A] (p) {};
  \node[bpmn task, rounded corners=0pt, right=of p] (B) {$B$};
  \node[bpmn event, right=of B] (o) {};
  \draw[bpmn flow] (i) -- (A);
  \draw[bpmn flow] (A) -- (p);
  \draw[bpmn flow] (p) -- (B);
  \draw[bpmn flow] (B) -- (o);
\end{tikzpicture}
\end{center}

\begin{examplebox}{Pattern: deferred choice (implicit XOR)}
Instead of inserting a routing transition, $i$ directly feeds both $A$ and $B$: the two transitions compete on the single token in $i$. Whichever fires first consumes the token, disabling the other. The branches reconverge at an XOR-join into $o$.
\[
  i \to \{A, B\}; \quad A \to p_A \to \text{XOR-join} \to o; \quad B \to p_B \to \text{XOR-join} \to o.
\]
The choice is now \textbf{implicit}: no routing transition ``decides''; the decision falls out of the interleaving semantics and of which environment event triggers $A$ or $B$ first.
\end{examplebox}

\begin{notebox}{Explicit choice $\neq$ implicit choice}
The two choice patterns are not interchangeable: the \emph{decision is made at different points in time}.

\begin{itemize}
  \item Explicit choice (XOR-split + XOR-join). The XOR-split transition $T$ \textbf{chooses} upfront, before $A$ or $B$ can start. In BPMN this maps to the \textbf{XOR gateway} ($\times$ diamond).
  \item Deferred choice. $A$ and $B$ are both pre-enabled; whichever fires first \emph{is} the choice. ``A or B chooses''. In BPMN this maps to the \textbf{event-based gateway} (pentagon-in-diamond).
\end{itemize}

The two constructs yield identical \emph{traces} but different \emph{behavioural} properties: with external triggers on $A$ and $B$, the deferred-choice variant lets the environment decide.
\end{notebox}

\begin{examplebox}{Pattern: one-or-more iteration (XOR-join + XOR-split)}
To execute $A$ \emph{at least once}, and possibly several times, the loop sandwiches $A$ between a join place and a split place. The output place $p_{\mathrm{loop}}^{\mathrm{end}}$ offers an XOR choice: \emph{either} fire $t_{\mathrm{repeat}}$ (a back-arc into $p_{\mathrm{loop}}^{\mathrm{start}}$, re-enabling $A$), \emph{or} fire $t_{\mathrm{exit}}$ into $o$.
\[
  i \to t_{\mathrm{in}} \to p_{\mathrm{loop}}^{\mathrm{start}} \to A \to p_{\mathrm{loop}}^{\mathrm{end}} \to \{t_{\mathrm{repeat}}, t_{\mathrm{exit}}\}; \quad t_{\mathrm{repeat}} \to p_{\mathrm{loop}}^{\mathrm{start}}; \quad t_{\mathrm{exit}} \to o.
\]
The XOR-join at $p_{\mathrm{loop}}^{\mathrm{start}}$ accepts the token from either $t_{\mathrm{in}}$ or $t_{\mathrm{repeat}}$. The XOR-split at $p_{\mathrm{loop}}^{\mathrm{end}}$ chooses whether to loop or to exit. The exit transition can fire only \emph{after} $A$ has fired at least once, so the iteration is guaranteed to execute at least one round.

\medskip
% Verified 2026-07-10: places i, p_loop_start, p_loop_end, o; transitions t_in, A,
% t_repeat, t_exit. Arcs i->t_in, t_in->p_loop_start, p_loop_start->A, A->p_loop_end,
% p_loop_end->t_repeat, p_loop_end->t_exit, t_repeat->p_loop_start, t_exit->o.
% Replay from M0={i:1}: t_in -> {p_loop_start}; A -> {p_loop_end}; now t_repeat and
% t_exit both enabled (XOR-split). t_exit -> {o} exits; t_repeat -> {p_loop_start}
% loops back to A. Only A feeds p_loop_end, so t_exit reaches o only after A fires
% >=1 time: one-or-more holds. Coords chosen to fit A4, no resizebox.
\begin{center}
\begin{tikzpicture}
  \node[bpmn event, minimum size=1.5em] (i) at (0,0) {};
  \node[bpmn task, rounded corners=0pt, minimum size=1.5em] (tin) at (1.4,0) {};
  \node[bpmn event, minimum size=1.5em] (ps) at (2.8,0) {};
  \node[bpmn task, rounded corners=0pt, minimum size=1.8em] (A) at (4.2,0) {};
  \draw[accent, thick] ($(A.north west)+(2pt,-2pt)$) rectangle ($(A.south east)+(-2pt,2pt)$);
  \node[bpmn event, minimum size=1.5em] (pe) at (5.6,0) {};
  \node[bpmn task, rounded corners=0pt, minimum size=1.5em] (trep) at (4.2,1.5) {};
  \node[bpmn task, rounded corners=0pt, minimum size=1.5em] (tex) at (7,0) {};
  \node[bpmn event, minimum size=1.5em] (o) at (8.4,0) {};
  \fill[accent] (i.center) circle (2pt);
  \draw[bpmn flow] (i) -- (tin);
  \draw[bpmn flow] (tin) -- (ps);
  \draw[bpmn flow] (ps) -- (A);
  \draw[bpmn flow] (A) -- (pe);
  \draw[bpmn flow] (pe) -- (tex);
  \draw[bpmn flow] (tex) -- (o);
  \draw[bpmn flow] (pe) |- (trep);
  \draw[bpmn flow] (trep) -| (ps);
  \node[below=1pt of i, font=\scriptsize] {$i$};
  \node[below=1pt of tin, font=\scriptsize] {$t_{\mathrm{in}}$};
  \node[below=1pt of ps, font=\scriptsize] {$p_{\mathrm{loop}}^{\mathrm{start}}$};
  \node[above=1pt of A, font=\scriptsize] {$A$};
  \node[below=1pt of pe, font=\scriptsize] {$p_{\mathrm{loop}}^{\mathrm{end}}$};
  \node[above=1pt of trep, font=\scriptsize] {$t_{\mathrm{repeat}}$};
  \node[below=1pt of tex, font=\scriptsize] {$t_{\mathrm{exit}}$};
  \node[below=1pt of o, font=\scriptsize] {$o$};
\end{tikzpicture}
\end{center}
\end{examplebox}


\begin{examplebox}{Pattern: zero-or-more iteration}
To allow \emph{zero or more} executions of $A$ (possibly none), $A$ sits on the \emph{back edge} of the loop, not on the forward path:
\[
  i \to t_{\mathrm{in}} \to p_2; \quad p_2 \to t \to p_3; \quad p_3 \to A \to p_2; \quad p_3 \to t_{\mathrm{exit}} \to o.
\]
When $p_3$ holds a token, both $A$ and $t_{\mathrm{exit}}$ are enabled. The token can fire $t_{\mathrm{exit}}$ immediately (zero executions of $A$), or cycle through $A$ and $t$ any number of times before exiting. The XOR is \emph{deferred}: $A$ and $t_{\mathrm{exit}}$ compete on the single token in $p_3$.

\medskip
\begin{center}
\begin{tikzpicture}
  \node[bpmn event, minimum size=1.5em] (i) at (0,0) {};
  \node[bpmn task, rounded corners=0pt, minimum size=1.5em] (tin) at (1.4,0) {};
  \node[bpmn event, minimum size=1.5em] (p2) at (2.8,0) {};
  \node[bpmn task, rounded corners=0pt, minimum size=1.5em] (t) at (4.2,0) {};
  \node[bpmn event, minimum size=1.5em] (p3) at (5.6,0) {};
  \node[bpmn task, rounded corners=0pt, minimum size=1.8em] (A) at (4.2,1.5) {};
  \draw[accent, thick] ($(A.north west)+(2pt,-2pt)$) rectangle ($(A.south east)+(-2pt,2pt)$);
  \node[bpmn task, rounded corners=0pt, minimum size=1.5em] (tex) at (7,0) {};
  \node[bpmn event, minimum size=1.5em] (o) at (8.4,0) {};
  \draw[bpmn flow] (i) -- (tin);
  \draw[bpmn flow] (tin) -- (p2);
  \draw[bpmn flow] (p2) -- (t);
  \draw[bpmn flow] (t) -- (p3);
  \draw[bpmn flow] (p3) -- (A);
  \draw[bpmn flow] (A) -- (p2);
  \draw[bpmn flow] (p3) -- (tex);
  \draw[bpmn flow] (tex) -- (o);
  \node[below=1pt of i, font=\scriptsize] {$i$};
  \node[below=1pt of tin, font=\scriptsize] {$t_{\mathrm{in}}$};
  \node[below=1pt of p2, font=\scriptsize] {$p_2$};
  \node[below=1pt of t, font=\scriptsize] {$t$};
  \node[below=1pt of p3, font=\scriptsize] {$p_3$};
  \node[above=1pt of A, font=\scriptsize] {$A$};
  \node[below=1pt of tex, font=\scriptsize] {$t_{\mathrm{exit}}$};
  \node[below=1pt of o, font=\scriptsize] {$o$};
\end{tikzpicture}
\end{center}
The double square marks $A$ as a sub-process (see the Sub-processes subsection). The same glyph recurs in all the pattern figures below.
\end{examplebox}

\begin{examplebox}{Zero-or-more iteration: sugared}
The decorations make the routing explicit: the \emph{inbound} transition is an XOR-join, accepting the token from $i$ or from the loop-back place, and the \emph{outbound} one an XOR-split, routing it to $o$ or back through $A$.

\medskip
\begin{center}
\begin{tikzpicture}
  \node[bpmn event, minimum size=1.5em] (i) at (0,0) {};
  \node[bpmn task, rounded corners=0pt, minimum size=1.5em] (tin) at (1.6,0) {};
  \fill[rulegray] (tin.north west) -- (tin.south west) -- ($(tin.west)+(5pt,0)$) -- cycle;
  \node[bpmn event, minimum size=1.5em] (pm) at (3.6,0) {};
  \node[bpmn task, rounded corners=0pt, minimum size=1.5em] (tout) at (5.6,0) {};
  \fill[rulegray] (tout.north east) -- (tout.south east) -- ($(tout.east)+(-5pt,0)$) -- cycle;
  \node[bpmn event, minimum size=1.5em] (o) at (7.2,0) {};
  \node[bpmn event, minimum size=1.5em] (pr) at (5.6,1.5) {};
  \node[bpmn task, rounded corners=0pt, minimum size=1.8em] (A) at (3.6,1.5) {};
  \draw[accent, thick] ($(A.north west)+(2pt,-2pt)$) rectangle ($(A.south east)+(-2pt,2pt)$);
  \node[bpmn event, minimum size=1.5em] (pl) at (1.6,1.5) {};
  \draw[bpmn flow] (i) -- (tin);
  \draw[bpmn flow] (tin) -- (pm);
  \draw[bpmn flow] (pm) -- (tout);
  \draw[bpmn flow] (tout) -- (o);
  \draw[bpmn flow] (tout) -- (pr);
  \draw[bpmn flow] (pr) -- (A);
  \draw[bpmn flow] (A) -- (pl);
  \draw[bpmn flow] (pl) -- (tin);
  \node[below=1pt of i, font=\scriptsize] {$i$};
  \node[above=1pt of A, font=\scriptsize] {$A$};
  \node[below=1pt of o, font=\scriptsize] {$o$};
\end{tikzpicture}
\end{center}
\end{examplebox}

\begin{examplebox}{Capacity constraint: one serve per time}
When many activations may reach the sub-process $A$, but $A$ must serve them \emph{one at a time}, the standard pattern uses a binary semaphore: a shared place $s$ holding exactly one token (the ``capacity''). The place is wired to the transitions \emph{bracketing} $A$, not to $A$ itself: the entry transition $t_{\mathrm{in}}$ consumes the capacity token together with the activation, and the exit transition $t_{\mathrm{out}}$ returns it when $A$'s run completes. While $A$ runs the capacity is held, so no second activation can enter. Bracketing matters because $A$ is a sub-process, hence non-atomic; a plain transition fires atomically and could never be observed ``busy''.

\medskip
\begin{center}
\begin{tikzpicture}
  \node[bpmn event, minimum size=1.5em] (i) at (0,0) {};
  \node[bpmn task, rounded corners=0pt, minimum size=1.5em] (tin) at (1.5,0) {};
  \node[bpmn event, minimum size=1.5em] (p1) at (3,0) {};
  \node[bpmn task, rounded corners=0pt, minimum size=1.8em] (A) at (4.5,0) {};
  \draw[accent, thick] ($(A.north west)+(2pt,-2pt)$) rectangle ($(A.south east)+(-2pt,2pt)$);
  \node[bpmn event, minimum size=1.5em] (p2) at (6,0) {};
  \node[bpmn task, rounded corners=0pt, minimum size=1.5em] (tout) at (7.5,0) {};
  \node[bpmn event, minimum size=1.5em] (o) at (9,0) {};
  \node[bpmn event, minimum size=1.5em] (s) at (4.5,1.7) {};
  \fill[accent] (s.center) circle (2pt);
  \draw[bpmn flow] (i) -- (tin);
  \draw[bpmn flow] (tin) -- (p1);
  \draw[bpmn flow] (p1) -- (A);
  \draw[bpmn flow] (A) -- (p2);
  \draw[bpmn flow] (p2) -- (tout);
  \draw[bpmn flow] (tout) -- (o);
  \draw[bpmn flow] (s) -- (tin);
  \draw[bpmn flow] (tout) -- (s);
  \node[below=1pt of tin, font=\scriptsize] {$t_{\mathrm{in}}$};
  \node[below=1pt of A, font=\scriptsize] {$A$};
  \node[below=1pt of tout, font=\scriptsize] {$t_{\mathrm{out}}$};
  \node[above=1pt of s, font=\scriptsize] {$s$};
\end{tikzpicture}
\end{center}
\end{examplebox}

\begin{examplebox}{Capacity constraint: mutual exclusion}
The same idea generalises to two sub-processes $A$ and $B$ that must not execute concurrently. Each runs in its own chain, and a shared \emph{mutex} place $m$ holding one token sits between them, wired crosswise to the bracketing transitions: the entry transition of each chain consumes the mutex token, the exit transition returns it. While one chain holds the token, the other chain's entry is disabled.

\medskip
\begin{center}
\begin{tikzpicture}
  \node[bpmn event, minimum size=1.5em] (ia) at (0,0) {};
  \node[bpmn task, rounded corners=0pt, minimum size=1.5em] (tina) at (1.5,0) {};
  \node[bpmn event, minimum size=1.5em] (pa1) at (3,0) {};
  \node[bpmn task, rounded corners=0pt, minimum size=1.8em] (A) at (4.5,0) {};
  \draw[accent, thick] ($(A.north west)+(2pt,-2pt)$) rectangle ($(A.south east)+(-2pt,2pt)$);
  \node[bpmn event, minimum size=1.5em] (pa2) at (6,0) {};
  \node[bpmn task, rounded corners=0pt, minimum size=1.5em] (touta) at (7.5,0) {};
  \node[bpmn event, minimum size=1.5em] (oa) at (9,0) {};
  \node[bpmn event, minimum size=1.5em] (ib) at (0,-3) {};
  \node[bpmn task, rounded corners=0pt, minimum size=1.5em] (tinb) at (1.5,-3) {};
  \node[bpmn event, minimum size=1.5em] (pb1) at (3,-3) {};
  \node[bpmn task, rounded corners=0pt, minimum size=1.8em] (B) at (4.5,-3) {};
  \draw[accent, thick] ($(B.north west)+(2pt,-2pt)$) rectangle ($(B.south east)+(-2pt,2pt)$);
  \node[bpmn event, minimum size=1.5em] (pb2) at (6,-3) {};
  \node[bpmn task, rounded corners=0pt, minimum size=1.5em] (toutb) at (7.5,-3) {};
  \node[bpmn event, minimum size=1.5em] (ob) at (9,-3) {};
  \node[bpmn event, minimum size=1.5em] (m) at (4.5,-1.5) {};
  \fill[accent] (m.center) circle (2pt);
  \draw[bpmn flow] (ia) -- (tina);
  \draw[bpmn flow] (tina) -- (pa1);
  \draw[bpmn flow] (pa1) -- (A);
  \draw[bpmn flow] (A) -- (pa2);
  \draw[bpmn flow] (pa2) -- (touta);
  \draw[bpmn flow] (touta) -- (oa);
  \draw[bpmn flow] (ib) -- (tinb);
  \draw[bpmn flow] (tinb) -- (pb1);
  \draw[bpmn flow] (pb1) -- (B);
  \draw[bpmn flow] (B) -- (pb2);
  \draw[bpmn flow] (pb2) -- (toutb);
  \draw[bpmn flow] (toutb) -- (ob);
  \draw[bpmn flow] (m) -- (tina);
  \draw[bpmn flow] (m) -- (tinb);
  \draw[bpmn flow] (touta) -- (m);
  \draw[bpmn flow] (toutb) -- (m);
  \node[above=1pt of A, font=\scriptsize] {$A$};
  \node[below=1pt of B, font=\scriptsize] {$B$};
  \node[right=2pt of m, font=\scriptsize] {$m$};
\end{tikzpicture}
\end{center}
\end{examplebox}

\begin{examplebox}{Capacity constraint: alternation}
$A$ and $B$ execute one time each, $A$ first. The wiring follows the mutual-exclusion pattern, but with \emph{two} control places: $d$ (initially marked) enables the entry of the $A$-chain; the exit of the $A$-chain produces into $c$, which enables the entry of the $B$-chain; the exit of the $B$-chain produces back into $d$. With the control token starting in $d$, the $A$-chain must complete before the $B$-chain can start.

\medskip
\begin{center}
\begin{tikzpicture}
  \node[bpmn event, minimum size=1.5em] (ia) at (0,0) {};
  \node[bpmn task, rounded corners=0pt, minimum size=1.5em] (tina) at (1.5,0) {};
  \node[bpmn event, minimum size=1.5em] (pa1) at (3,0) {};
  \node[bpmn task, rounded corners=0pt, minimum size=1.8em] (A) at (4.5,0) {};
  \draw[accent, thick] ($(A.north west)+(2pt,-2pt)$) rectangle ($(A.south east)+(-2pt,2pt)$);
  \node[bpmn event, minimum size=1.5em] (pa2) at (6,0) {};
  \node[bpmn task, rounded corners=0pt, minimum size=1.5em] (touta) at (7.5,0) {};
  \node[bpmn event, minimum size=1.5em] (oa) at (9,0) {};
  \node[bpmn event, minimum size=1.5em] (ib) at (0,-3) {};
  \node[bpmn task, rounded corners=0pt, minimum size=1.5em] (tinb) at (1.5,-3) {};
  \node[bpmn event, minimum size=1.5em] (pb1) at (3,-3) {};
  \node[bpmn task, rounded corners=0pt, minimum size=1.8em] (B) at (4.5,-3) {};
  \draw[accent, thick] ($(B.north west)+(2pt,-2pt)$) rectangle ($(B.south east)+(-2pt,2pt)$);
  \node[bpmn event, minimum size=1.5em] (pb2) at (6,-3) {};
  \node[bpmn task, rounded corners=0pt, minimum size=1.5em] (toutb) at (7.5,-3) {};
  \node[bpmn event, minimum size=1.5em] (ob) at (9,-3) {};
  \node[bpmn event, minimum size=1.5em] (c) at (3.4,-1.5) {};
  \node[bpmn event, minimum size=1.5em] (d) at (5.6,-1.5) {};
  \fill[accent] (d.center) circle (2pt);
  \draw[bpmn flow] (ia) -- (tina);
  \draw[bpmn flow] (tina) -- (pa1);
  \draw[bpmn flow] (pa1) -- (A);
  \draw[bpmn flow] (A) -- (pa2);
  \draw[bpmn flow] (pa2) -- (touta);
  \draw[bpmn flow] (touta) -- (oa);
  \draw[bpmn flow] (ib) -- (tinb);
  \draw[bpmn flow] (tinb) -- (pb1);
  \draw[bpmn flow] (pb1) -- (B);
  \draw[bpmn flow] (B) -- (pb2);
  \draw[bpmn flow] (pb2) -- (toutb);
  \draw[bpmn flow] (toutb) -- (ob);
  \draw[bpmn flow] (d) -- (tina);
  \draw[bpmn flow] (touta) -- (c);
  \draw[bpmn flow] (c) -- (tinb);
  \draw[bpmn flow] (toutb) -- (d);
  \node[above=1pt of A, font=\scriptsize] {$A$};
  \node[below=1pt of B, font=\scriptsize] {$B$};
  \node[left=2pt of c, font=\scriptsize] {$c$};
  \node[right=2pt of d, font=\scriptsize] {$d$};
\end{tikzpicture}
\end{center}
\end{examplebox}

\begin{examplebox}{Question time: which model is more reasonable?}
Two candidate workflow nets for the same travel-booking process (\textsc{travel\_request} $\to$ \textsc{book\_hotel} and \textsc{book\_flight} $\to$ \textsc{send\_plan}).

\begin{itemize}
  \item \textbf{First model.} Place $p_2$ feeds both \textsc{book\_hotel} and \textsc{book\_flight}, and both produce into the shared place $p_3$, an \emph{implicit} (deferred) choice: the two transitions compete on the single token, so exactly \emph{one} booking happens per case.
  \item \textbf{Second model.} \textsc{travel\_request} is an AND-split producing into two parallel places; the bookings run concurrently and their output places synchronise on \textsc{send\_plan}. \emph{Both} bookings happen per case.
\end{itemize}

\medskip
\begin{center}
\begin{tikzpicture}
  \node[bpmn event, minimum size=1.5em] (p1) at (0,0) {};
  \fill[accent] (p1.center) circle (2pt);
  \node[bpmn task, rounded corners=0pt, font=\scriptsize\sffamily] (tr) at (1.7,0) {travel request};
  \node[bpmn event, minimum size=1.5em] (p2) at (3.4,0) {};
  \node[bpmn task, rounded corners=0pt, font=\scriptsize\sffamily] (bh) at (5.1,0.9) {book hotel};
  \node[bpmn task, rounded corners=0pt, font=\scriptsize\sffamily] (bf) at (5.1,-0.9) {book flight};
  \node[bpmn event, minimum size=1.5em] (p3) at (6.8,0) {};
  \node[bpmn task, rounded corners=0pt, font=\scriptsize\sffamily] (sp) at (8.4,0) {send plan};
  \node[bpmn event, minimum size=1.5em] (p4) at (10,0) {};
  \draw[bpmn flow] (p1) -- (tr);
  \draw[bpmn flow] (tr) -- (p2);
  \draw[bpmn flow] (p2) -- (bh);
  \draw[bpmn flow] (p2) -- (bf);
  \draw[bpmn flow] (bh) -- (p3);
  \draw[bpmn flow] (bf) -- (p3);
  \draw[bpmn flow] (p3) -- (sp);
  \draw[bpmn flow] (sp) -- (p4);
  \node[below=1pt of p1, font=\scriptsize] {$p_1$};
  \node[below=1pt of p2, font=\scriptsize] {$p_2$};
  \node[below=1pt of p3, font=\scriptsize] {$p_3$};
  \node[below=1pt of p4, font=\scriptsize] {$p_4$};
\end{tikzpicture}

\medskip
\begin{tikzpicture}
  \node[bpmn event, minimum size=1.5em] (p1) at (0,0) {};
  \fill[accent] (p1.center) circle (2pt);
  \node[bpmn task, rounded corners=0pt, font=\scriptsize\sffamily] (tr) at (1.7,0) {travel request};
  \node[bpmn event, minimum size=1.5em] (p2) at (3.4,0.9) {};
  \node[bpmn event, minimum size=1.5em] (p3) at (3.4,-0.9) {};
  \node[bpmn task, rounded corners=0pt, font=\scriptsize\sffamily] (bh) at (5.1,0.9) {book hotel};
  \node[bpmn task, rounded corners=0pt, font=\scriptsize\sffamily] (bf) at (5.1,-0.9) {book flight};
  \node[bpmn event, minimum size=1.5em] (p4) at (6.8,0.9) {};
  \node[bpmn event, minimum size=1.5em] (p5) at (6.8,-0.9) {};
  \node[bpmn task, rounded corners=0pt, font=\scriptsize\sffamily] (sp) at (8.4,0) {send plan};
  \node[bpmn event, minimum size=1.5em] (p6) at (10,0) {};
  \draw[bpmn flow] (p1) -- (tr);
  \draw[bpmn flow] (tr) -- (p2);
  \draw[bpmn flow] (tr) -- (p3);
  \draw[bpmn flow] (p2) -- (bh);
  \draw[bpmn flow] (p3) -- (bf);
  \draw[bpmn flow] (bh) -- (p4);
  \draw[bpmn flow] (bf) -- (p5);
  \draw[bpmn flow] (p4) -- (sp);
  \draw[bpmn flow] (p5) -- (sp);
  \draw[bpmn flow] (sp) -- (p6);
  \node[below=1pt of p1, font=\scriptsize] {$p_1$};
  \node[above=1pt of p2, font=\scriptsize] {$p_2$};
  \node[below=1pt of p3, font=\scriptsize] {$p_3$};
  \node[above=1pt of p4, font=\scriptsize] {$p_4$};
  \node[below=1pt of p5, font=\scriptsize] {$p_5$};
  \node[below=1pt of p6, font=\scriptsize] {$p_6$};
\end{tikzpicture}
\end{center}

Which one is more reasonable? The second. A travel booking needs both the hotel and the flight reserved before the plan can be sent: the two tasks are independent and complementary, not mutually exclusive. The first model would book only one of the two and ship the case as ``planned''.
\end{examplebox}

\subsection{Triggers}

\begin{definitionbox}{Triggers}
Execution constraints can depend on the \emph{environment} in which a process is enacted. In a workflow net, transitions can be decorated with information on \textbf{who} or \textbf{what} is responsible for firing them; such annotations are called \textbf{triggers}.
\end{definitionbox}

A trigger is one of three kinds: a \textbf{user} takes the initiative, an \textbf{external} event arrives (typically a message), or a \textbf{timer} elapses. A transition with no trigger fires \emph{automatically} as soon as it is enabled. WoPeD's transition-properties dialog therefore offers four options, counting the automatic default as one of them.\footnote{\url{http://woped.dhbw-karlsruhe.de/woped/}} Each carries a small icon on top of the transition box:

\medskip

\begin{center}
\begin{tabular}{cl}
\toprule
\textbf{Symbol} & \textbf{Meaning} \\
\midrule
(empty)        & \textbf{Automatic}: no trigger, the task fires on its own. \\
$\downarrow$   & \textbf{User}: a human user takes initiative. \\
$\boxtimes$    & \textbf{External}: external event required (e.g.\ message). \\
$\bigcirc$\,\textit{clock} & \textbf{Timer}: activity starts when timer elapses. \\
\bottomrule
\end{tabular}
\end{center}

The shapes match the BPMN event-type catalogue.\footnote{Weske M., \emph{Business Process Management}, Springer, 2007.}

\begin{examplebox}{Trigger example: request handling}
A small workflow net illustrates the three non-automatic triggers in action:

\begin{itemize}
  \item \textsc{send\_request} carries a \emph{User} trigger: a clerk pushes the request out;
  \item \textsc{collect\_response} carries an \emph{External} trigger: it fires only when an incoming message arrives from the counterpart;
  \item \textsc{send\_reminder} carries a \emph{Timer} trigger and loops back to the shared waiting place: if no response arrives within a timeout, a reminder is sent and the wait restarts;
  \item \textsc{archive\_file} is undecorated: it fires automatically once the response has been collected.
\end{itemize}

\medskip
\begin{center}
\begin{tikzpicture}
  \node[bpmn event, minimum size=1.5em] (i) at (0,0) {};
  \node[bpmn task, rounded corners=0pt, font=\scriptsize\sffamily] (sr) at (1.8,0) {send request};
  \node[above=1pt of sr, font=\scriptsize, text=accent] {User};
  \node[bpmn event, minimum size=1.5em] (pw) at (3.8,0) {};
  \node[bpmn task, rounded corners=0pt, font=\scriptsize\sffamily] (cr) at (5.8,0) {collect response};
  \node[above=1pt of cr, font=\scriptsize, text=accent] {External};
  \node[bpmn event, minimum size=1.5em] (p2) at (7.9,0) {};
  \node[bpmn task, rounded corners=0pt, font=\scriptsize\sffamily] (af) at (9.6,0) {archive file};
  \node[bpmn event, minimum size=1.5em] (o) at (11.3,0) {};
  \node[bpmn task, rounded corners=0pt, font=\scriptsize\sffamily] (rem) at (3.8,-1.4) {send reminder};
  \node[right=1pt of rem, font=\scriptsize, text=accent] {Timer};
  \draw[bpmn flow] (i) -- (sr);
  \draw[bpmn flow] (sr) -- (pw);
  \draw[bpmn flow] (pw) -- (cr);
  \draw[bpmn flow] (cr) -- (p2);
  \draw[bpmn flow] (p2) -- (af);
  \draw[bpmn flow] (af) -- (o);
  \draw[bpmn flow] (pw) to[bend right=30] (rem);
  \draw[bpmn flow] (rem) to[bend right=30] (pw);
  \node[below=1pt of i, font=\scriptsize] {$i$};
  \node[below=1pt of o, font=\scriptsize] {$o$};
\end{tikzpicture}
\end{center}
\end{examplebox}

\begin{examplebox}{Realistic trigger example: credit-card order handling}
A larger workflow net for credit-card order handling: nine places ($i$, $A,\dots,G$, $o$) and nine transitions, several of them decorated. \textsc{pack complete or incomplete}, \textsc{ship}, \textsc{receive} and \textsc{backorder} carry User triggers; \textsc{cancel order} carries a Timer trigger (the order is cancelled if not packed within the time-out); \textsc{update billing info updated or cancelled} carries an External trigger (it waits for the bank's message); \textsc{start}, \textsc{spam customer} and the XOR-split \textsc{charge credit card success or failure} fire automatically.

\medskip
\begin{center}
\resizebox{\linewidth}{!}{%
\begin{tikzpicture}
  \node[bpmn event, minimum size=1.5em] (i) at (0,0) {};
  \fill[accent] (i.center) circle (2pt);
  \node[bpmn task, rounded corners=0pt, font=\scriptsize\sffamily] (start) at (2.9,0) {start};
  \node[bpmn event, minimum size=1.5em] (A) at (2.9,-0.8) {};
  \node[bpmn task, rounded corners=0pt, font=\scriptsize\sffamily] (charge) at (6.7,-0.8) {charge credit card};
  \fill[rulegray] (charge.north east) -- (charge.south east) -- ($(charge.east)+(-5pt,0)$) -- cycle;
  \node[bpmn event, minimum size=1.5em] (D) at (10.7,-0.8) {};
  \node[bpmn task, rounded corners=0pt, font=\scriptsize\sffamily] (pack) at (14.2,-0.8) {pack};
  \fill[rulegray] (pack.north east) -- (pack.south east) -- ($(pack.east)+(-5pt,0)$) -- cycle;
  \node[above=1pt of pack, font=\scriptsize, text=accent] {User};
  \node[bpmn event, minimum size=1.5em] (G) at (16.8,-0.8) {};
  \node[bpmn task, rounded corners=0pt, font=\scriptsize\sffamily] (ship) at (19.2,-0.8) {ship};
  \node[above=1pt of ship, font=\scriptsize, text=accent] {User};
  \node[bpmn event, minimum size=1.5em] (B) at (6.7,-1.6) {};
  \node[bpmn task, rounded corners=0pt, font=\scriptsize\sffamily] (recv) at (10.7,-1.6) {receive};
  \node[right=1pt of recv, font=\scriptsize, text=accent] {User};
  \node[bpmn event, minimum size=1.5em] (E) at (14.2,-1.6) {};
  \node[bpmn task, rounded corners=0pt, font=\scriptsize\sffamily] (spam) at (6.7,-2.4) {spam customer};
  \node[bpmn event, minimum size=1.5em] (F) at (10.7,-2.4) {};
  \node[bpmn task, rounded corners=0pt, font=\scriptsize\sffamily] (back) at (14.2,-2.4) {backorder};
  \node[right=1pt of back, font=\scriptsize, text=accent] {User};
  \node[bpmn event, minimum size=1.5em] (C) at (8.5,-3.2) {};
  \node[bpmn task, rounded corners=0pt, font=\scriptsize\sffamily] (cancel) at (11.8,-3.2) {cancel order};
  \node[above=1pt of cancel, font=\scriptsize, text=accent] {Timer};
  \node[bpmn task, rounded corners=0pt, font=\scriptsize\sffamily] (upd) at (4.8,-4.0) {update billing info};
  \fill[rulegray] (upd.north east) -- (upd.south east) -- ($(upd.east)+(-5pt,0)$) -- cycle;
  \node[left=1pt of upd, font=\scriptsize, text=accent] {External};
  \node[bpmn event, minimum size=1.5em] (o) at (19.2,-4.0) {};
  \draw[bpmn flow] (i) -- (start);
  \draw[bpmn flow] (start) -- (A);
  \draw[bpmn flow] (A) -- (charge);
  \draw[bpmn flow] (charge) -- (D);
  \draw[bpmn flow] (charge) -- (B);
  \draw[bpmn flow] (D) -- (pack);
  \draw[bpmn flow] (pack) -- (G);
  \draw[bpmn flow] (pack) -- (E);
  \draw[bpmn flow] (G) -- (ship);
  \draw[bpmn flow] (ship) -- (o);
  \draw[bpmn flow] (E) -- (back);
  \draw[bpmn flow] (back) -- (F);
  \draw[bpmn flow] (F) -- (recv);
  \draw[bpmn flow] (recv) -- (D);
  \draw[bpmn flow] (B) -- (spam);
  \draw[bpmn flow] (spam) -- (C);
  \draw[bpmn flow] (C) -- (cancel);
  \draw[bpmn flow] (cancel) -- (o);
  \draw[bpmn flow] (C) -- (upd);
  \draw[bpmn flow] (upd) -- (A);
  \draw[bpmn flow] (upd) -- (o);
  \node[below=1pt of i, font=\scriptsize] {$i$};
  \node[left=2pt of A, font=\scriptsize] {$A$};
  \node[above=1pt of D, font=\scriptsize] {$D$};
  \node[above=1pt of G, font=\scriptsize] {$G$};
  \node[left=2pt of B, font=\scriptsize] {$B$};
  \node[above=1pt of E, font=\scriptsize] {$E$};
  \node[below=1pt of F, font=\scriptsize] {$F$};
  \node[below=1pt of C, font=\scriptsize] {$C$};
  \node[below=1pt of o, font=\scriptsize] {$o$};
\end{tikzpicture}}
\end{center}
The hatched right halves mark the three XOR-splits (\textsc{charge}: success routes to $D$, failure to $B$; \textsc{pack}: complete to $G$, incomplete to $E$; \textsc{update billing info}: updated back to $A$ for a new charge, cancelled to $o$).
\end{examplebox}

\begin{notebox}{Explicit vs implicit choice: with triggers}
Triggers make the explicit/implicit contrast (drawn above) operationally sharp. Under an explicit XOR-split a routing transition fires automatically and enables only the chosen branch, so $A$ and $B$ are never concurrently enabled. Under a deferred choice both are enabled and whichever triggers first (a user click on $A$, a timer expiring on $B$) consumes the shared token and disables the other. The BPMN mapping is the one already given: XOR gateway for the explicit case, event-based gateway for the deferred one.
\end{notebox}

\begin{examplebox}{Exercise: Car Damage, net design}
An insurance company uses the following procedure for the processing of claims:

\begin{itemize}
  \item every claim, reported by a customer, is \textbf{registered};
  \item after the registration, the claim is \textbf{classified} into one of two categories: \textbf{simple} or \textbf{complex};
  \item for \emph{simple} claims, two tasks must be executed: \textsc{check\_insurance} and \textsc{phone\_garage}; these tasks are \emph{independent} of each other;
  \item for \emph{complex} claims, three tasks are required: \textsc{check\_insurance}, \textsc{check\_damage\_history}, and \textsc{phone\_garage}; these must be executed \emph{sequentially} in the order specified;
  \item after executing the two-or-three tasks, a \textbf{decision} is taken with two outcomes: \textbf{OK} (positive) or \textbf{NOK} (negative);
  \item if the decision is positive, the company will \textbf{pay};
  \item in any event, the company \textbf{sends a letter} to the customer.
\end{itemize}

Design the workflow net.

\medskip
\textbf{Solution.} The target workflow net combines all the patterns of the chapter.

\begin{itemize}
  \item \textbf{Register} (mandatory): a single transition right after the source place \textsc{begin}, producing into the control place $c_1$.
  \item \textbf{Classify} as XOR-split: from $c_1$ the token is routed either into the place \textsc{simple} or into the place \textsc{complex}.
  \item \textbf{Simple branch}: an AND-split from \textsc{simple} into $\{c_2, c_3\}$ feeds the two independent tasks \textsc{phone\_garage} (into $c_5$) and \textsc{check\_insurance} (into $c_6$); an AND-join synchronises $\{c_5, c_6\}$ into $c_8$.
  \item \textbf{Complex branch}: a strict sequence \textsc{complex} $\to$ \textsc{check\_insurance} $\to c_4 \to$ \textsc{check\_history} $\to c_7 \to$ \textsc{phone\_garage} $\to c_8$. The place $c_8$ acts as the XOR-join of the two branches.
  \item \textbf{Decide} as XOR-split from $c_8$: the token goes either to \textsc{OK} or to \textsc{NOK}.
  \item \textbf{Pay} on the OK branch only: \textsc{OK} $\to$ \textsc{pay} $\to c_9$; the XOR-join \textsc{send\_letter} accepts the token from $c_9$ or from \textsc{NOK}, so a letter is sent in any event.
  \item \textbf{End}: \textsc{send\_letter} produces into the sink place \textsc{end}. All human tasks carry User triggers; \textsc{pay} and the AND-split/join fire automatically.
\end{itemize}

\medskip
\begin{center}
\resizebox{0.95\linewidth}{!}{%
\begin{tikzpicture}
  \node[bpmn event, minimum size=1.5em] (begin) at (0,0) {};
  \fill[accent] (begin.center) circle (2pt);
  \node[bpmn task, rounded corners=0pt, font=\scriptsize\sffamily] (reg) at (0,-1.2) {register};
  \node[left=1pt of reg, font=\scriptsize, text=accent] {User};
  \node[bpmn event, minimum size=1.5em] (c1) at (0,-2.4) {};
  \node[bpmn task, rounded corners=0pt, font=\scriptsize\sffamily] (cls) at (1.7,-2.4) {classify};
  \fill[rulegray] (cls.north east) -- (cls.south east) -- ($(cls.east)+(-5pt,0)$) -- cycle;
  \node[above=1pt of cls, font=\scriptsize, text=accent] {User};
  \node[bpmn event, minimum size=1.5em] (simple) at (3.4,-1.2) {};
  \node[bpmn task, rounded corners=0pt, font=\scriptsize\sffamily] (asplit) at (4.9,-1.2) {AND split};
  \node[bpmn event, minimum size=1.5em] (c2) at (6.4,-0.2) {};
  \node[bpmn task, rounded corners=0pt, font=\scriptsize\sffamily] (pg1) at (8,-0.2) {phone garage};
  \node[above=1pt of pg1, font=\scriptsize, text=accent] {User};
  \node[bpmn event, minimum size=1.5em] (c5) at (9.6,-0.2) {};
  \node[bpmn event, minimum size=1.5em] (c3) at (6.4,-2.2) {};
  \node[bpmn task, rounded corners=0pt, font=\scriptsize\sffamily] (ci1) at (8,-2.2) {check insurance};
  \node[above=1pt of ci1, font=\scriptsize, text=accent] {User};
  \node[bpmn event, minimum size=1.5em] (c6) at (9.6,-2.2) {};
  \node[bpmn task, rounded corners=0pt, font=\scriptsize\sffamily] (ajoin) at (11.1,-1.2) {AND join};
  \node[bpmn event, minimum size=1.5em] (c8) at (12.6,-1.2) {};
  \node[bpmn event, minimum size=1.5em] (complex) at (3.4,-3.6) {};
  \node[bpmn task, rounded corners=0pt, font=\scriptsize\sffamily] (ci2) at (5.1,-3.6) {check insurance};
  \node[below=1pt of ci2, font=\scriptsize, text=accent] {User};
  \node[bpmn event, minimum size=1.5em] (c4) at (6.7,-3.6) {};
  \node[bpmn task, rounded corners=0pt, font=\scriptsize\sffamily] (ch) at (8.2,-3.6) {check history};
  \node[below=1pt of ch, font=\scriptsize, text=accent] {User};
  \node[bpmn event, minimum size=1.5em] (c7) at (9.7,-3.6) {};
  \node[bpmn task, rounded corners=0pt, font=\scriptsize\sffamily] (pg2) at (11.2,-3.6) {phone garage};
  \node[below=1pt of pg2, font=\scriptsize, text=accent] {User};
  \node[bpmn task, rounded corners=0pt, font=\scriptsize\sffamily] (dec) at (1,-4.9) {decide};
  \fill[rulegray] (dec.north east) -- (dec.south east) -- ($(dec.east)+(-5pt,0)$) -- cycle;
  \node[left=1pt of dec, font=\scriptsize, text=accent] {User};
  \node[bpmn event, minimum size=1.5em] (ok) at (2.8,-4.6) {};
  \node[bpmn task, rounded corners=0pt, font=\scriptsize\sffamily] (pay) at (4.4,-4.6) {pay};
  \node[bpmn event, minimum size=1.5em] (c9) at (6,-4.6) {};
  \node[bpmn event, minimum size=1.5em] (nok) at (2.8,-5.8) {};
  \node[bpmn task, rounded corners=0pt, font=\scriptsize\sffamily] (sl) at (6,-5.8) {send letter};
  \fill[rulegray] (sl.north west) -- (sl.south west) -- ($(sl.west)+(5pt,0)$) -- cycle;
  \node[below=1pt of sl, font=\scriptsize, text=accent] {User};
  \node[bpmn event, minimum size=1.5em] (end) at (7.8,-5.8) {};
  \draw[bpmn flow] (begin) -- (reg);
  \draw[bpmn flow] (reg) -- (c1);
  \draw[bpmn flow] (c1) -- (cls);
  \draw[bpmn flow] (cls) -- (simple);
  \draw[bpmn flow] (cls) -- (complex);
  \draw[bpmn flow] (simple) -- (asplit);
  \draw[bpmn flow] (asplit) -- (c2);
  \draw[bpmn flow] (asplit) -- (c3);
  \draw[bpmn flow] (c2) -- (pg1);
  \draw[bpmn flow] (pg1) -- (c5);
  \draw[bpmn flow] (c3) -- (ci1);
  \draw[bpmn flow] (ci1) -- (c6);
  \draw[bpmn flow] (c5) -- (ajoin);
  \draw[bpmn flow] (c6) -- (ajoin);
  \draw[bpmn flow] (ajoin) -- (c8);
  \draw[bpmn flow] (complex) -- (ci2);
  \draw[bpmn flow] (ci2) -- (c4);
  \draw[bpmn flow] (c4) -- (ch);
  \draw[bpmn flow] (ch) -- (c7);
  \draw[bpmn flow] (c7) -- (pg2);
  \draw[bpmn flow] (pg2) -- (c8);
  \draw[bpmn flow] (c8.south) -- (12.6,-4.2) -- (1,-4.2) -- (dec.north);
  \draw[bpmn flow] (dec) -- (ok);
  \draw[bpmn flow] (dec) -- (nok);
  \draw[bpmn flow] (ok) -- (pay);
  \draw[bpmn flow] (pay) -- (c9);
  \draw[bpmn flow] (c9) -- (sl);
  \draw[bpmn flow] (nok) -- (sl);
  \draw[bpmn flow] (sl) -- (end);
  \node[right=2pt of begin, font=\scriptsize] {begin};
  \node[left=2pt of c1, font=\scriptsize] {$c_1$};
  \node[above=1pt of simple, font=\scriptsize] {simple};
  \node[above=1pt of c2, font=\scriptsize] {$c_2$};
  \node[above=1pt of c5, font=\scriptsize] {$c_5$};
  \node[below=1pt of c3, font=\scriptsize] {$c_3$};
  \node[below=1pt of c6, font=\scriptsize] {$c_6$};
  \node[above=1pt of c8, font=\scriptsize] {$c_8$};
  \node[above=1pt of complex, font=\scriptsize] {complex};
  \node[below=1pt of c4, font=\scriptsize] {$c_4$};
  \node[below=1pt of c7, font=\scriptsize] {$c_7$};
  \node[above=1pt of ok, font=\scriptsize] {OK};
  \node[above=1pt of c9, font=\scriptsize] {$c_9$};
  \node[below=1pt of nok, font=\scriptsize] {NOK};
  \node[below=1pt of end, font=\scriptsize] {end};
\end{tikzpicture}}
\end{center}
\end{examplebox}
